\section{Proof of the main theorem}
\label{sec:conclusion}

We now assemble the exponent construction and interpolation results. The
reshaping parameters will depend on the point count $r$, but a single
threshold will allow every sufficiently large integer $r$. The final
passage from jets to intersection numbers uses ordinary curve
multiplicities, including at singular points of the curves.

\subsection{The exponent sets for every sufficiently large point count}

Fix the polarized variety $(X,L)$ of Theorem~\ref{thm:main}, and write
$V=L^n$. Choose an integer $D>0$ for which $DL$ is very ample.
Proposition~\ref{prop:flag-simplex} supplies additive support sets $S_k$ with
\[
 S_k\subset kP_0\cap\Z^n,
 \qquad
 \#S_k=h^0(X,kL).
\]
Here $P_0=\Delta(1/D,\ldots,1/D,D^{n-1}V)$.
Let $\ell>0$ be the smallest intercept of $P_0$. Choose an integer $r_0$
such that
\begin{equation}
 \label{eq:threshold}
 (V/r)^{1/n}<\ell/100\qquad(r\ge r_0).
\end{equation}
Such a choice exists because $X$ and $L$ are fixed.

Fix any integer $r\ge r_0$, and put $w=(V/r)^{1/n}$. Apply
Proposition~\ref{prop:reshape} to the first two intercepts, replacing them
by $w$ and their product divided by $w$. Keep the intercept $w$ fixed, and
pair the other new intercept with the next untouched one. Continue until
$n-1$ intercepts equal $w$.

Every pair of input intercepts $A,B$ in this procedure satisfies
$A,B\ge\ell$. Indeed, an untouched intercept has this property, and the
new carried intercept $AB/w$ is greater than $A$, since $B\ge\ell>w$.
Thus the hypothesis of Proposition~\ref{prop:reshape} holds at every step:
\[
 w\max(3/A,2/B)\le\frac{3w}{\ell}<\frac1{16}.
\]
The product of all intercepts remains $V$. We obtain finite sets $F_k$ in
\begin{equation}
 \label{eq:final-simplex}
 kP_*,\qquad
 P_* =\Delta(w,\ldots,w,V/w^{n-1})
     =\Delta(w,\ldots,w,rw),
\end{equation}
with
\[
 F_k+F_l\subset F_{k+l},
 \qquad
 \#F_k=h^0(X,kL)=\frac{V}{n!}k^n+o(k^n).
\]
Since $\vol(P_*)=V/n!$, these sets satisfy the hypotheses of
Proposition~\ref{prop:asymptotic-jets}.

For each $0<\theta<w$ and all sufficiently large $k$, that proposition
gives simultaneous jets through degree $\lfloor k\theta\rfloor$ for
$\M(F_k)$ at an $r$-tuple of distinct torus points. Each application of
Proposition~\ref{prop:reshape} transfers this surjectivity to the preceding
support set. Proposition~\ref{prop:flag-simplex} finally transfers it to
$H^0(X,kL)$ at some tuple of distinct points of $X$. In particular, for
every such $\theta$ and $k$, the set
\begin{equation}
 \label{eq:surjectivity-locus}
 \left\{\mathbf p\in U_r(X):
 H^0(X,kL)\longrightarrow
 \bigoplus_{i=1}^r\Jet_{p_i}^{\lfloor k\theta\rfloor}(kL)
 \text{ is surjective}\right\}
\end{equation}
is nonempty. The tuple witnessing this assertion may depend on both
$\theta$ and $k$.

\subsection{A common very-general locus}

For fixed integers $k\ge1$ and $m\ge0$, the locus of tuples at which
$H^0(X,kL)$ surjects onto jets through degree $m$ is Zariski open in
$U_r(X)$. Indeed, the bundle of principal parts of order $m$ of $kL$ is
locally free because $X$ is smooth. Pulling it back along each coordinate
projection of $U_r(X)$ gives an algebraic vector-bundle map
\[
 H^0(X,kL)\otimes\mathcal O_{U_r(X)}
 \longrightarrow
 \bigoplus_{i=1}^r\operatorname{pr}_i^*\mathcal P^m(kL).
\]
Its surjectivity locus is defined by the nonvanishing of maximal minors.

Choose $0<\theta_j<w$ with $\theta_j\uparrow w$. For each $j$, choose
$k_j$ so that the locus in \eqref{eq:surjectivity-locus} is nonempty and
$\lfloor k\theta_j\rfloor\ge1$ for every $k\ge k_j$. Denote this open
locus by $O_{j,k}$ and set
\begin{equation}
 \label{eq:exceptional-set}
 Z_r=\bigcup_{j\ge1}\ \bigcup_{k\ge k_j}
       \bigl(U_r(X)\setminus O_{j,k}\bigr).
\end{equation}
Every set in this countable union is a proper Zariski-closed subset of
$U_r(X)$.

The complement is nonempty. The complex points of $U_r(X)$ form a
nonempty locally compact Hausdorff space, and hence a Baire space. Since
$U_r(X)$ is smooth and irreducible, every proper algebraic closed subset
has empty interior in its complex topology. A countable union of these
closed subsets cannot cover $U_r(X)$. Thus every tuple outside $Z_r$
simultaneously has all the jet-surjectivity properties indexed by $j$ and
$k\ge k_j$.

\subsection{From jets to ordinary curve multiplicities}

The relation between Seshadri constants and jet separation appears in
Demailly's one-point formulation \cite[Section~6]{Demailly1992} and in the
weighted multipoint framework of \citet[Section~5]{Ito2013}. The following elementary
implication is the one needed here. We give its local proof, including at
singular points of the curves.

\begin{lemma}[Jets bound curve intersections]
\label{lem:jets-curves}
Let $X$ be a smooth complex projective variety, let $A$ be a line bundle
on $X$, and let $p_1,\ldots,p_r$ be distinct points. Suppose that, for an
integer $m\ge1$, the map
\[
 H^0(X,A)\longrightarrow\bigoplus_{i=1}^r\Jet_{p_i}^{m}(A)
\]
is surjective. Then every integral curve $C\subset X$ satisfies
\[
 A\cdot C\ge m\sum_{i=1}^r\mult_{p_i}C
\]
provided that $C$ meets at least one of the points.
\end{lemma}

\begin{proof}
Choose a point $p_i\in C$. The local ring $R=\mathcal O_{C,p_i}$ is a
one-dimensional Noetherian local domain. Its maximal ideal $\frakm_R$
satisfies
\[
 \frakm_R^m/\frakm_R^{m+1}\ne0.
\]
Otherwise Nakayama's lemma would give $\frakm_R^m=0$, contrary to
$\dim R=1$. The quotient map from the ambient local ring to $R$ lifts a
nonzero element of this quotient to a jet on $X$ with no terms of degree
less than $m$. Prescribe this jet at $p_i$ and the zero jet at every other
marked point. By surjectivity, a section $s\in H^0(X,A)$ has these jets.
It vanishes to order at least $m$ at every marked point, and its restriction
to $C$ is nonzero.

We recall the local intersection bound for this restriction. In a
one-dimensional Noetherian local domain $(R,\frakm_R)$, let
$0\ne f\in\frakm_R^m$. For every integer $q\ge1$, multiplication by
powers of $f$ gives
\[
 \length(R/(f^q))=q\,\length(R/(f)).
\]
Since $(f^q)\subset\frakm_R^{mq}$, the Hilbert--Samuel formula yields
\[
 q\,\length(R/(f))
 \ge\length(R/\frakm_R^{mq})
 =mq\,e(R)+O(1),
\]
where $e(R)$ is the ordinary multiplicity. Dividing by $q$ and letting
$q\to\infty$ gives
\begin{equation}
 \label{eq:local-curve-bound}
 \length(R/(f))\ge m\,e(R).
\end{equation}

Apply \eqref{eq:local-curve-bound} to the local equations of $s|_C$ at
the marked points on $C$. The zero scheme of this nonzero section is an
effective Cartier divisor on the integral projective curve, and its total
length is $\deg(A|_C)=A\cdot C$. Summing the local contributions proves
the assertion.
\end{proof}

Fix a tuple $\mathbf p\in U_r(X)\setminus Z_r$. Applying
Lemma~\ref{lem:jets-curves} with $A=kL$ and
$m=\lfloor k\theta_j\rfloor$ gives
\[
 \eps(X,L;\mathbf p)\ge\frac{\lfloor k\theta_j\rfloor}{k}
 \qquad(k\ge k_j).
\]
First let $k\to\infty$ for each fixed $j$, and then let $j\to\infty$.
We obtain
\begin{equation}
 \label{eq:main-lower-bound}
 \eps(X,L;\mathbf p)\ge w=(V/r)^{1/n}.
\end{equation}

\subsection{The volume bound}

For completeness, we derive the reverse inequality for every tuple of
distinct points. Let $\pi:Y\to X$ be the blow-up of such a tuple, with
exceptional divisors $E_1,\ldots,E_r$. For $t\ge0$, put
\[
 D_t=\pi^*L-t\sum_{i=1}^r E_i.
\]
If an integral curve $\Gamma$ is contained in $E_i\simeq\mathbb P^{n-1}$,
then $D_t|_{E_i}=t\mathcal O_{\mathbb P^{n-1}}(1)$ is nonnegative on
$\Gamma$. Every other integral curve is the strict transform of an
integral curve $C\subset X$, and
\begin{equation}
 \label{eq:strict-transform-intersection}
 D_t\cdot\Gamma=L\cdot C-t\sum_{i=1}^r\mult_{p_i}C.
\end{equation}
Here $E_i\cdot\Gamma=\mult_{p_i}C$ also when $C$ is singular. If
$p_i\notin C$, both sides vanish. For $p_i\in C$, identify the strict
transform locally with the blow-up of $\mathcal O_{C,p_i}$ at its maximal
ideal. Its intersection
with $E_i$ is the zero-dimensional exceptional fiber
$\operatorname{Proj}(\operatorname{gr}_{\frakm}\mathcal O_{C,p_i})$,
whose length is the Hilbert--Samuel multiplicity. This is the degree of
the effective Cartier divisor cut out by $E_i$ on the strict transform.

It follows from \eqref{eq:strict-transform-intersection} that $D_t$ is
nef whenever $0\le t<\eps(X,L;\mathbf p)$. A nef real divisor has
nonnegative top self-intersection. The exceptional divisors are disjoint,
$\pi^*L|_{E_i}$ is trivial, and $E_i^n=(-1)^{n-1}$. Therefore
\[
 0\le D_t^n=V-rt^n.
\]
This inequality for all $0\le t<\eps(X,L;\mathbf p)$ proves the universal
bound
\begin{equation}
 \label{eq:volume-upper-bound}
 \eps(X,L;\mathbf p)\le(V/r)^{1/n}.
\end{equation}
Together, \eqref{eq:main-lower-bound} and \eqref{eq:volume-upper-bound}
give the equality asserted in Theorem~\ref{thm:main} at every tuple outside
$Z_r$. Equation~\eqref{eq:strict-transform-intersection} also shows that
$D_w$ is nef there, and its top self-intersection is zero. The threshold
\eqref{eq:threshold} was chosen before fixing $r$, so this proves the
theorem for every integer $r\ge r_0$.
